\chapter{Statistical Methods for the Phantom and Clinical Comparisons}
\label{app:90y_statistical_models}

These are the full specifications for the comparisons in Chapter~\ref{chapter:paper}. In the phantom analysis, $c$ identifies the reconstruction method, $v$ the sphere, and $\mathrm{bs}$ the bootstrap replicate. In the clinical models, $p$ identifies the patient and $l$ the lesion. The reference method is \wdtnv. The symbols $\beta_0$ and $\alpha_c$ denote the intercept and method effects within each model; parameters are estimated separately for each endpoint. Random effects and residual errors are assumed mutually independent, with the Gaussian distributions specified below.

\section{Phantom Recovery Coefficients}

For each sphere $v$ and method $c\neq c_0$, inference for RC uses the paired log-ratios
\begin{equation}
d_{\mathrm{bs},c,v} := \log\!\left(\mathrm{RC}^{(c)}_{\mathrm{bs},v}\right)-\log\!\left(\mathrm{RC}^{(c_0)}_{\mathrm{bs},v}\right),
\qquad \mathrm{bs}=1,\ldots,N_{\mathrm{bs}}.
\end{equation}
The log-scale contrast is estimated by the mean paired log-ratio
\begin{equation}
\widehat{\Delta}_{c,v} = \frac{1}{N_{\mathrm{bs}}}\sum_{\mathrm{bs}=1}^{N_{\mathrm{bs}}} d_{\mathrm{bs},c,v},
\end{equation}
so that $\exp(\widehat{\Delta}_{c,v})$ is the ratio of the geometric-mean RC of method $c$ to that of \wdtnv. The 95\% confidence interval is the exponentiated Student-$t$ interval for the mean of $d_{\mathrm{bs},c,v}$ ($N_{\mathrm{bs}}-1$ degrees of freedom). Under the null hypothesis of exchangeability between $c$ and $c_0$, each $d_{\mathrm{bs},c,v}$ is equally likely to take either sign. Enumerating all $2^{N_{\mathrm{bs}}}$ sign patterns $\bm{\sigma}\in\{-1,1\}^{N_{\mathrm{bs}}}$ gives the exact two-sided p-value
\begin{equation}
p_{c,v} = \frac{1}{2^{N_{\mathrm{bs}}}}\sum_{\bm{\sigma}} \mathbf{1}\!\left(\left|\sum_{\mathrm{bs}=1}^{N_{\mathrm{bs}}}\sigma_{\mathrm{bs}}\, d_{\mathrm{bs},c,v}\right| \ge \left|\sum_{\mathrm{bs}=1}^{N_{\mathrm{bs}}} d_{\mathrm{bs},c,v}\right|\right).
\end{equation}

\section{Clinical Contrast and Background Noise}

We analyse $\log(\mathrm{TBR})$ using a linear mixed-effects model with fixed method effects and random intercepts for patient and lesion-within-patient:
\begin{equation}
\begin{aligned}
y_{p,l,c} &:= \log\!\left(\mathrm{TBR}_{p,l,c}\right)
= \beta_0 + \alpha_c + u_p + v_{p,l} + \varepsilon_{p,l,c},\\
u_p &\overset{\mathrm{iid}}{\sim} \mathcal{N}(0,\tau^2_{\mathrm{p}}),
\qquad
v_{p,l} \overset{\mathrm{iid}}{\sim} \mathcal{N}(0,\tau^2_{\mathrm{l}}),
\qquad
\varepsilon_{p,l,c} \overset{\mathrm{iid}}{\sim} \mathcal{N}(0,\tau^2).
\end{aligned}
\end{equation}
Background noise is analysed analogously at the patient level using
\begin{equation}
\begin{aligned}
y^{\mathrm{CoV}}_{p,c} &:= \log\!\left(\mathrm{CoV}_{p,c}(\%)\right)
= \beta_0 + \alpha_c + u_p + \varepsilon_{p,c},\\
u_p &\overset{\mathrm{iid}}{\sim} \mathcal{N}(0,\tau^2_{\mathrm{p,CoV}}),
\qquad
\varepsilon_{p,c} \overset{\mathrm{iid}}{\sim} \mathcal{N}(0,\psi^2).
\end{aligned}
\end{equation}
